\paragraph{Dressed charge profile and current response.}
A positive-energy Pauli source with a compact Maxwell field retains the
magnetic commutator, both spin blocks and every integer electric flux.
The electric and magnetic coefficients are $1/32$ and $1/8$, with a positive
covariant Wilson coefficient $1/16$. Physical probes use the native charge
$Q=-1$ and retain its opposite quantum reference. A complete charge-position
measurement gives $R_x^2=\langle x^2\rangle$, independently of the current
spectrum. On the four-site plaquette the exact elastic form factor is
$F(k)=1-\langle x^2\rangle(1-\cos k)$. Padded rational residuals certify the
infinite-source ground profile. The regulated electron has
$R_x\simeq0.12118$ cell and maximum first-Brillouin-interval suppression
$1-F\simeq0.02937$. The core is therefore below the spatial resolution needed
to demonstrate an order-one onset. The zero-Wilson controls retain the
smaller heavy-family responses.

The physical dipole $D=Qx$ supplies $J=i[H,D]$ and the required contact
$T=-[D,[D,H]]$. The exact static Ward sum is
$\langle T\rangle=2\sum_n(E_n-E_0)|D_{n0}|^2$.
The contact-inclusive imaginary-frequency response is
\[
 \Gamma(\nu)=2\nu^2\left\langle Dg,
 \frac{H-E_0}{(H-E_0)^2+\nu^2}Dg\right\rangle.
\]
Exact padded resolvent residuals and ground-state bounds give absolute
errors below $10^{-9}$ on four complete infinite-source response values.
For the regulated electron these are approximately
$2.96653\times10^{-5}$, $4.42359\times10^{-4}$,
$3.52069\times10^{-3}$ and $6.74284\times10^{-3}$ at
$\nu=1/64,1/16,1/4,1$, respectively, in electron-energy units.
All individual spectral weights and complete probe outcomes are retained.
The source is an intraband construction; these results supply neither
pair-producing vacuum running nor a physical electron-core fit.
